\scanpage{35}
\mspage{13}
Si $i,j\in I^*$, il suffit \inserted{(grâce à f))}
de vérifier que pour tt
$\beta$, on a
\[
X_i\cap X_j\cap T_\beta=\bigcup_{k\le i,j}X_k\cap T_\beta
\]
(il suffit de prouver $\subset$, car $\supset$ trivial). Si
$X_i\cap X_j\cap T_\beta=\varnothing$ c'est trivial, on peut donc supposer
$X_i\cap X_j\cap T_\beta\ne\varnothing$
\deleted{\unsure{i.e.} $X_i\cap X_j\cap$}
donc $X_i\cap T_\beta\ne\varnothing$, $X_j\cap T_\beta\ne\varnothing$
i.e. $i,j\in I_\beta$. Alors \deleted{\unsure{de}} cela résulte de
(f) induit sur $\unsure{\overline T_\beta|\dot T_\beta}=\dot X_\beta\subset X^*$,
et de la formule (19). Donc

\uline{Thm} Les systèmes (17), satisfaisant aux
conditions \deleted{\unsure{tt}} (14) des stratifications, et aux
conditions de transversalité explicitées plus
haut, équivalent aux systèmes (18), satisfaisant
aux conditions (20), en plus des conditions
e) f) exprimant que les $(X_i^*)_{i\in I^*}$ forment une
stratification de $X^*$.

8. \uline{Structure cônique sur un espace stratifié}

Pour $X$ stratifié de type $I$ par des
$(X_i)_{i\in I}$, avec des sous-espaces bordants donnés
\inserted{$(B_\alpha)_{\alpha\in\Lambda}$}
et des germes de \inserted{voisinages} cylindriques des $B_\alpha$, tels que
\deleted{\unsure{ces bords} \unsure{de} \unsure{manière}}
les $(B_\alpha)$ soient mutuellement transverses, et
transverses aux $X_i$, on va \deleted{construire} définir
\margin{34}
